\documentclass[12pt]{article}
\usepackage[T1]{fontenc}
\usepackage[utf8]{inputenc}
\usepackage[brazil]{babel}
\usepackage[margin=1in]{geometry}
\setlength{\parindent}{0pt}
\begin{document}
\section*{Tema 39}
Processo: PEDILEF 0506910-51.2005.4.05.8013/ AL\\
Situacao: Julgado\\
Ramo: DIREITO PREVIDENCIÁRIO\\
Questao: Saber se o trabalhador que preencheu requisito de carência do benefício de aposentadoria por idade, e que na data do óbito, não mais detinha qualidade de segurado e tampouco havia implementado idade mínima, pode ser instituidor de benefício de pensão por morte.\\
Tese: Para a concessão de pensão por morte de rurícola é necessário que o instituidor tenha, na data do óbito, a qualidade de segurado ou tenha implementado, antes de falecer, todos os requisitos para a concessão de aposentadoria por idade rural, tanto a carência quanto a idade mínima.\\
Fonte: https://www.cjf.jus.br/phpdoc/virtus/
\end{document}
